
\documentclass{BeamerTemplate}
\title{Module 5 - Lois conjointes pour variables discrètes}
\subtitle{STT-1000 Probabilités et statistique}
\date{\today}

\usepackage{adjustbox}


\begin{document}
\AtBeginSection{\frame{\sectionpage}}

\begin{frame}

\titlepage

\end{frame}

%-------------------------------------------------------------------------------------------------------------------------------------





\section{Vecteurs aléatoires}


\begin{frame}[t]{Vecteurs aléatoires}

Dans bien des situations, on peut s'intéresser au comportement aléatoire de plusieurs variables \alerter{simultanément}.
\vspace{0.1 in}

\begin{itemize}
 \item Le réservoir d'un barrage reçoit l'eau de 4 rivières.
 On sera intéressé par le comportement \alerter{conjoint} de $\mathbf{X}=(X,
 Y,X_3,X_4)$, où $X_i$ est le débit de la $i$-ème rivière. 
 \item Un système en parallèle fonctionne tant qu'au moins
 un de ses trois composants fonctionne. Soit $Y_i$, la variable qui
 indique si le $i$-ème composant fonctionne. On sera intéressé
 par le comportement \alerter{conjoint} de $\mathbf{Y}=(Y_1,Y_2,Y_3)$.
\end{itemize}

\vspace{0.1 in}

\textbf{Pour notre cours: uniquement le cas de 2 variables aléatoires discrètes.
}

\end{frame}


\section{Loi conjointe de 2 variables aléatoires}

\begin{frame}{Fonction de probabilité conjointe}

 \begin{defn}{}
  La \alerter{fonction de probabilité conjointe} de deux v.a.
  discrètes $X$ et $Y$ est la fonction $p(x,y)$ qui associe à toute
  valeur possible $(x,y)$ de $(X,Y)$ la probabilité
  $P(X=x\mbox{\ et\ }Y=y)$:
  $$p(x,y)=P(X=x\mbox{\ et\ }Y=Y).$$
 \end{defn}
 
 \vspace{0.1in}
 Puisque c'est une fonction de probabilité, on a :
 \begin{itemize}
 \item $0\leq p(x,y) \leq 1$ pour tout $(x,y)$
 \item $\sum\limits_{x}\sum\limits_{y}p(x,y)=1$
 \end{itemize}
 

\end{frame}

\begin{frame}{Exemple}

On considère un réseau d'alimentation électrique divisé en deux parties: le réseau $\mathcal{X}$ alimente $2$ maisons et le réseau $\mathcal{Y}$ alimente $3$ maisons. Parfois, il y a des problèmes avec le réseau et le courant ne se rend pas à toutes les maisons. \\
%C'est un peu petit, mais bon ça pourrait être le village en dessous du sapin de Noel...

\vspace{0.1 in}
On définit les variables suivantes: \\
$X$ = nombre de maisons du réseau $\mathcal{X}$ sans électricité \\
$Y$ = nombre de maisons du réseau $\mathcal{Y}$ sans électricité \\

\vspace{0.1in}
Une fonction de probabilité conjointe possible pour ces variables aléatoires est:
\begin{center}
\begin{tabular}{cc|cccc}
$p(x,y)$ & & \multicolumn{3}{c}{$y$} \\
 & & 0 & 1 & 2 & 3 \\ \cline{1-6}
 & 0 & 0.2 & 0.15  & 0.05 & 0 \\
$x$ & 1 & 0.1 & 0.1 & 0.05 & 0.1\\
 & 2 & 0.1 & 0.05 & 0.05 & 0.05 \\
\end{tabular}
\end{center}

\end{frame}

\section{Lois marginales}




\begin{frame}[t]{Lois marginales}

 \begin{defn}{}
  Soit $p(x,y)$, la fonction de probabilité conjointe de
  $X$ et $Y$. Alors les \alerter{fonctions de probabilité marginales}
  de $X$ et $Y$ sont respectivement données par
  \begin{align*}
   p_{X}(x)&=P(X=x)=\sum_{\alerter{y}}p(x,y)\\
   p_Y(y)&=P(Y=y)=\sum_{\alerter{x}}p(x,y).
  \end{align*}
 \end{defn}



\end{frame}

\begin{frame}[t]{Retour à l'exemple}

On peut calculer les deux fonctions de probabilité marginales:


\end{frame}

\begin{frame}[t]{Retour à l'exemple}

Quelle est la probabilité qu'il n'y ait aucune maison sans électricité dans le réseau $\mathcal{X}$?





\end{frame}

\section{Lois conditionnelles}



\begin{frame}[t]{Lois conditionnelles}


 \begin{defn}{}
  Soit $p(x,y)$, la fonction de probabilité conjointe de
  $X$ et $Y$. Alors les \alerter{fonctions de probabilité conditionnelles}
  de $X$ sachant $Y$ et de $Y$ sachant $X$ sont respectivement données par
  \begin{align*}
   p_{X|y}(x)&=P(X=x|Y=y)=\frac{P(X=x\text{\ et\ }Y=y)}{P(Y=y)}
   =\frac{p(x,y)}{p_Y(y)}\\
   & \mbox{(fonction de $x$ pour une valeur de $y$ fixe)} \\[15pt]
   p_{Y|x}(y)&=P(Y=y|X=x)=\frac{P(X=x\text{\ et\ }Y=y)}{P(X=x)}
   =\frac{p(x,y)}{p_X(x)}\\
     & \mbox{(fonction de $y$ pour une valeur de $x$ fixe)}
  \end{align*}
 \end{defn}
 


\end{frame}

\begin{frame}[t]{Retour à l'exemple }

On peut calculer les fonctions de probabilité conditionnelles:

\end{frame}

\begin{frame}[t]{Retour à l'exemple}

Quelle est la probabilité qu'il y ait une maison du réseau $\mathcal{Y}$ sans électricité s'il y a 2 maisons du réseau $\mathcal{X}$ sans électricité?


\end{frame}

\begin{frame}[t]{Espérance conditionnelle}

{\small
 Puisque la fonction de probabilité conditionnelle d'une v.a. est une fonction
 de probabilité en bonne et due forme, on peut s'en servir pour calculer des espérances,
 variances, etc.

 \begin{defn}{}
   Les \alerter{espérances conditionnelles} de $X$ sachant $Y$ et de $Y$ sachant
   $X$ sont respectivement données par
   \begin{align*}
     E[X|y]&=\sum\limits_{x}xp_{X|y}(x)\\
     E[Y|x]&=\sum\limits_{y}yp_{Y|x}(y).
   \end{align*}
 \end{defn}
 
Nous avons donc également
\begin{itemize}
\item $E[g(X)|y]=\sum\limits_{x}g(x)p_{X|y}(x)$
\item $E[g(Y)|x]=\sum\limits_{y}g(y)p_{Y|x}(y)$
\end{itemize}

}

\end{frame}

\begin{frame}[t]{Retour à l'exemple }
Quel est l'espérance du nombre de maisons sans électricité dans le réseau $\mathcal{Y}$ s'il y a 2 maisons sans électricité dans le réseau $\mathcal{X}$? Quelle est sa variance?

\end{frame}


\section{Indépendance de variables aléatoires}

\begin{frame}{Événements indépendants}
Rappel:
 \begin{defn}{}
  Soient $A$ et $B$ deux événements de $S$. On dit que $A$ et $B$
  sont \alerter{indépendants}, dénoté $A\perp B$, si, et seulement si,
  $$P(A\cap B) = P(A)P(B).$$
 \end{defn}
 
 

   Supposons que $P(B)>0$, de sorte que $P(A|B)$ est bien définie.

   \begin{itemize} \itemsep0.1in
    \item Si $A\perp B$, alors $P(A|B)=P(A)$.
    \item Si $P(A)>0$ aussi, alors $P(A|B)=P(A)\Leftrightarrow A\perp B$.
   \end{itemize}
   


\end{frame}


\begin{frame}[t]{Retour à l'exemple}

{\small
Est-ce que les événements suivants sont indépendants: \\
$E_1$ = Il y a deux maisons sans électricité dans le réseau $\mathcal{X}$. \\
$E_2$ = Il y a au moins une maison sans électricité dans le réseau $\mathcal{Y}$?
}


\end{frame}


\begin{frame}[t]{Indépendance de variables aléatoires}

{\small

 \begin{defn}{}
   Deux v.a. $X$ et $Y$ sont \alerter{indépendantes} si, et seulement si,
   $$p(x,y)=p_X(x)p_Y(y), \forall x, y$$
 \end{defn}
 \vspace{0.2in}
 \begin{itemize}
  \item Si on a $p(x,y)$ et que l'on veut savoir si $X\perp Y$, on doit obtenir $p_X(x)$ et $p_Y(y)$ et ensuite on
  vérifie si $p(x,y)=p_X(x)p_Y(y)$ pour toutes les paires de valeurs possibles.
  \item Si on a $p_X(x)$ et $p_Y(y)$ et que l'on nous demande de calculer
  $p(x,y)$, il n'est pas possible de le faire ... sauf si on nous
  dit que $X\perp Y$, auquel cas $p(x,y)=p_X(x)p_Y(y)$.
 \end{itemize}
}
\end{frame}

\begin{frame}{Conséquence de l'indépendance}


\begin{thm}{}
Deux v.a. $X$ et $Y$ sont indépendantes si, et seulement si
$$p_{Y|x}(y)=p_{Y}(y)\mbox{\ \ et\ \ }
p_{X|y}(x)=p_{X}(x)$$
pour toutes les valeurs possibles de $x$ et $y$.
\end{thm}

\end{frame}

\begin{frame}[t]{Retour à l'exemple}

Le nombre de maisons sans électricité dans le réseau $\mathcal{X}$ est-il indépendant du nombre de maisons sans électricité dans le réseau $\mathcal{Y}$?

\end{frame}

\begin{frame}[t]{Un autre exemple}
Soient $X$ et $Y$ deux v.a. ayant comme fonctions de
probabilité respectives $p_X(x)=1/3$ pour $x=1,~2$
et $p_X(x)=1/6$ pour $x=3,~4$, et $p_Y(y)=1/2$ pour
$y=4,~5$. Quelle est la fonction de probabilité conjointe
de $X$ et $Y$?

%Impossible de répondre avec l'information donnée!!!

%On peut le calculer si on suppose que $X$ et $Y$ sont indépendantes?}
\end{frame}

\section{La covariance et la corrélation}

\begin{frame}[t]{Covariance}
Si $X$ et $Y$ ne sont pas indépendantes, on peut mesurer leur association linéaire.

\begin{defn}{}
 La \alerter{covariance} entre deux v.a. $X$ et $Y$,
 dénotée $\sigma_{X,Y}$ ou $Cov(X,Y)$, est
 donnée par
 \begin{align*}
  \sigma_{XY}=Cov(X,Y)&=E[(X-E(X))(Y-E(Y))]
 \end{align*}
% où $E[g(X)h(Y)]=\sum\limits_{x}\sum\limits_{y}g(x)h(y)p(x,y)$.
\end{defn}

\vspace{0.2in}
On peut également écrire la covariance sous la forme suivante:

$$ Cov(X,Y) = E[XY]-E[X]E[Y]$$

La preuve est similaire à celle de $V(X) = E(X^2) - E(X)^2$.

\end{frame}

\begin{frame}[t]{Interprétation de la valeur de la covariance}
\begin{itemize}
 \item{\bf Association positive:} Une valeur $\sigma_{XY}>0$ signifie que les
 deux v.a. ont tendance à prendre des valeurs supérieures (ou inférieures) à leur moyenne
 respective \alerter{en même temps}.
 \item {\bf Association négative:} Une valeur $\sigma_{XY}<0$ signifie que
 $X$ a tendance à prendre des valeurs \alerter{supérieures} à sa moyenne quand
 $Y$ prend des valeurs \alerter{inférieures à sa moyenne} et vice-versa.
\end{itemize}

\vspace{0.1in}

\textbf{Unités:\\}
L'unité de mesure de $Cov(X,Y)$ est l'unité de mesure de $X$ multipliée par l'unité de mesure de $Y$. C'est donc assez difficile à interpréter...
\end{frame}

\begin{frame}[t]{Propriétés de la covariance}
 \begin{itemize}\itemsep4mm

 \item $Cov(X, X)=Var(X)$
 \item $Cov(aX, bY)=abCov(X, Y)$
 \item $Cov(X+c, Y+d)=Cov(X, Y)$
 \item $V(X+Y)=V(X)+V(Y)+2 Cov(X,Y)$
 \item $V(X-Y)=V(X)+V(Y)-2 Cov(X,Y)$

\item {\small $V(a_0+a_1X+\cdots+a_kX_k)=
{\color{blue}{ a_1^2\sigma^2_1+\cdots+a_k^2\sigma^2_k}}
+2\sum\limits_{i<i'}\sum
a_ia_{i'}Cov(X_i,X_{i'})$}
\end{itemize}
\end{frame}


\begin{frame}[t]%{Événements mutuellement exclusifs}

\centerline{\LARGE \bf \alert{QUIZ !}}
\begin{itemize}\itemsep1.5cm
\item  Que peut-on dire de $Cov(5,\, X)$? % 0
\item Vrai ou Faux: $Cov(X,\,3Y)=Cov(Y,\,3X)$? % Vrai
\item Si $X$ est mesurée en $^\circ C$ et $Y$ en secondes, quelles sont les unités de la covariance?
\end{itemize}

\end{frame}

\begin{frame}[t]{Coefficient de corrélation}

Mesure standardisée de la covariance:

\begin{defn}{}
 Le \alerter{coefficient de corrélation} entre deux variables aléatoires $X$ et $Y$ est donné par
 $$\rho=\frac{Cov(X,Y)}{\sqrt{V[X]V[Y]}}=\frac{\sigma_{XY}}{\sigma_X\sigma_Y}.$$
\end{defn}

\vspace{0.1 in}
\textbf{Propriétés:}
\begin{itemize}
 \item $-1\le \rho\le 1$  \\[12pt]
 \item  Si $X\perp Y$, alors $Cov(X,Y)=0$ et donc $\rho=0$.\\ (théorème 3.2, p-97)
\end{itemize}
\end{frame}

\begin{frame}[t]{Retour à l'exemple}
Calculer la \underline{covariance} et la {corrélation} entre le nombre de maisons sans électricité dans le réseau $\mathcal{X}$ et le nombre de maisons sans électricité dans le réseau $\mathcal{Y}$.


\end{frame}

\begin{frame}[t]{Retour à l'exemple (suite)}



\end{frame}

\begin{frame}[t]{Attention!}

Deux pièges à éviter au sujet de la corrélation:\\
\vspace{0.1 in}
\begin{itemize}
\item $\rho(X,Y) = 0 \centernot\implies X \perp Y$ \\[12pt]
\vspace{1.4in}
\item $\rho(X,Y)$ mesure la corrélation \textbf{linéaire}. Deux variables peuvent être reliées sans que  $\rho(X,Y)$ ne puisse le détecter.
\end{itemize}


\end{frame}

\section{Combinaisons linéaires}

\begin{frame}[t]{Deux variables - Espérance}

Supposons deux v.a. $X$ et $Y$. Soit $Y = a_1X + a_2Y$.

$\alerter{E(Y)}  =  E[ a_1X + a_2Y]$\\
... 






\end{frame}

\begin{frame}[t]{Deux variables - Variance}

\vspace{0.2 in}
{\footnotesize
\begin{align*}
V(Y) &=  E(Y^2) - (E(Y))^2\\
 & = E[ (a_1X + a_2Y )^2] - [E(a_1X + a_2Y)]^2\\ 
 & = E( a_1^2X^2 + 2a_1 a_2 X Y + a_2^2Y^2 ) + [a_1E(X) + a_2 E(Y)]^2  \\
 & = ...\\
 &= a_1^2 V(X) + 2 a_1 a_2 Cov(X, Y) + a_2^2 V(Y)
\end{align*}
}

$$ \alerter{V(a_1X + a_2 Y) = a_1^2 V(X)  + a_2^2 V(Y) + 2 a_1 a_2 Cov(X, Y) }$$



\end{frame}




\begin{frame}[t]{Généralisation pour plusieurs variables}
Soient $X,\hdots,X_n$, des variables aléatoires et $a_0,a_1,\hdots,a_n$, des constantes.\\
 Soit $Y=a_0+a_1X+\cdots+a_nX_n$. Alors
\begin{itemize}
 \item[] $$E[Y]=a_0+\sum_{i=1}^na_iE[X_i]$$ 
 \item[] $$V[Y] =\sum_{i=1}^na_i^2Var[X_i]+2\sum\limits_{j}\sum\limits_{i<j}a_ia_jCov(X_i,X_j)$$
\end{itemize}


Si $X,\hdots,X_n$ sont des v.a. \alerter{indépendantes}, alors $Cov(X_i,X_j)=0$ et  $$V[Y]=\sum_{i=1}^na_i^2V[X_i].$$

%\uncover<4-> {\alert{$\Rightarrow$ Le second point implique que $Var(X-Y)\neq Var(X)-Var(Y)$ mais bien $Var(X-Y)=Var(X)+Var(Y)-2Cov(X,Y)$.}}

\end{frame}

\begin{frame}[t]{Exemple}
Soient $X,Y,X_3$ les débits de 3 rivières.\\
Soit $Y$ la hausse d'un lac où ces rivières se jettent ou prennent leur source. On peut écrire $Y=0.6X+0.5Y-1.6X_3$. \\

\vspace{0.1in}
On sait que $E(X)=10$, $E(Y)=20$, $E(X_3)=10$,
$Var(X)=4$, $Var(Y)=16$, $Var(X_3)=1$,
$Cov(X,Y)=1.6$ et que $X_3$ est indépendante de $X$ et de $Y$. \\

\vspace{0.1in}
Calculez l'espérance et la variance de $Y$, la hausse du lac.
\end{frame}

\begin{frame}[t]{Exemple}

\end{frame}


\begin{frame}{Bibliographie}
 	Ce document est rédigé avec Beamer ainsi qu'une classe .cls fourni par Jérôme Soucy. Les diapositives sont adaptées des diapositives créées par Anne-Sophie Charest 
 pour le cours STT-1000 ainsi que celles créées par Thierry Duchesne et Emmanuelle Reny-Nolin pour le cours STT-1900.
 	\printbibliography
 	\vfill
 	\footnotesize Pour signaler une erreur dans ce document, veuillez écrire à l'adresse :~\href{mailto:julien.miron@mat.ulaval.ca}{julien.miron@mat.ulaval.ca}.
 \end{frame}

\end{document}

